\documentclass[czech,a4paper,11pt]{article}
%\documentclass[]{article}
\usepackage[pdftex]{graphicx}
\usepackage{graphicx}
\usepackage[czech]{babel}		
\usepackage[T1]{fontenc}	
\usepackage[utf8]{inputenc} 	
\newenvironment{bottompar}{\par\vspace*{\fill}}{\clearpage}
\usepackage{enumerate}
\usepackage{url}

\newcommand{\HRule}{\rule{\linewidth}{0.5mm}}

\begin{document}

\begin{titlepage}
\begin{center}

\includegraphics[width=0.15\textwidth]{./cvut-logo-bw.pdf}~\\[1cm]

\textsc{\LARGE České vysoké učení technické }\\[1.5cm]

\textsc{\Large MI-MKY úkol 3}\\[0.5cm]

% Title
\HRule \\[0.4cm]
{ \huge \bfseries Lámání DLP pomocí Babystep-Giantstep algoritmu \\[0.4cm] }

\HRule \\[1.5cm]

% Author and supervisor
\begin{minipage}{0.4\textwidth}
\begin{flushleft} \large
\emph{Autor:}\\
Tomáš \textsc{Sušánka}
\end{flushleft}
\end{minipage}

\vfill

% Bottom of the page
{\large \today}

\end{center}
\end{titlepage}


\newpage
\tableofcontents

\newpage
\section{Zadání a řešení}
Úkolem bylo implementovat algoritmus Babystep-Giantstep, který řeší problém diskrétního algoritmu.

Pro zadané $g(x)$ a mé soukromé $h(x)$ vyjde:

\begin{table}[h!]
  \begin{center}
	0x4658ca487
	$x^{34}+x^{30}+x^{29}+x^{26}+x^{24}+x^{23}+x^{19}+x^{18}+x^{15}+x^{13}+x^{10}+x^7+x^2+x+1$
  \end{center}
\end{table}


\section{Implementace}

Po dlouhé úvaze jsem zvolil pro implementaci jazyk Python. Začal jsem sice v jazyce C, ale došel jsem k názoru, že pro řazení (a zachování indexů) bude lepší využít některý z \uv{vyšších jazyků}.

Nejdříve jsem celý algoritmus implementoval pro celá čísla, abych si plně uvědomil, jak algoritmus funguje. Poté jsem implementoval násobení a sčítání v tělese $GF(2^{35})$\mbox{*}.

Násobení je implementováno pomocí algoritmu Double-and-Add. Pro vynásobení $a=ax \bmod {f(x)}$ je použito LFSR (Linear feedback shift register). Oba tyto postupy jsme se učili v předmětu MI-BHW, a proto mi přišlo vhodné je použít.

Sčítání je pouze jednoduché využití operátoru XOR.

\section{Problémy}
Problémů jsem měl až nečekaně hodně. Při implementaci násobení a práci s bity se člověk velmi lehce sekne.

Největší problém se kterým jsem se setkal byla nízká efektivita programu. Poté co jsem odladil násobení i hlavní část programu, jsem dal oba celky dohromady a výpočet mi trval kolem 15 minut. Dlouhou dobu jsem nemohl objevit, kde je zakopaný pes, stále jsem hledal chybu v implementaci násobení. Chyba samozřejmě byla úplně někde jinde.

Přestože jsem hodnoty v tabulce seřadil, hledal jsem v nich stále lineárně. Je však zřejmé, že v tomto případě je nasnadě využít binární půlení. Ze složitosti $O(n^2)$ jsem se náhle dostal na složitost $O(n\log{}n)$, což dělá dramatický rozdíl. Řešení po této opravě trvá cca 15 vteřin.

\section{Vylepšení programu}

Program by se nejspíše mohl dále vylepšovat. Příhodné by bylo zvlášť implementovat mocnění, které lze lehce vytvořit tzv. \emph{vkládáním nul mezi}. Stále je však potřeba číslo zmodulovat, takže se myslím, že zrychlení bude relativně malé. 

Další možnost je si druhé mocniny předpočítat. Opět se sice nevyhneme modulu, ale je pravděpodobné, že toto řešení bude nejrychlejší.

Protože jsem nevěřil, že by tato vylepšení přinesla markantní rozdíly, zajímalo mě jaký je rozdíl mezi jazykem Python a C, který jsem opustil.

\section{Benchmark}

Vytvořil jsem tedy velmi jednoduchý benchmark. Napsal jsem dva malé prográmky v C a Python, která vynásobí 100\,000\,000 čísel. Pro jednoduchost programy nenásobí v rámci tělesa GF, ale obyčejná desetinná čísla. Jak bylo možné očekávat, C je výrazně rychlejší:

\begin{table}[h]
	\centering
    \begin{tabular}{ll}
    C      & 0.84 s \\
    Python & 27.1 s \\
    \end{tabular}
\end{table}


Pokud bych algoritmus implementovat v čistém C, dá se očekávat výrazné zrychlení. Musíme však brát v potaz, že by autor musel implementovat řazení a hledání, případně použít C++. Tam už by rozdíl nebyl takový. Je tedy nutné brát benchmark s rezervou.


\section{Příloha}

Příloha obsahuje:

\begin{itemize}
	\item main.py a gf.py - řešení jako takové, hlavní soubor je main.py
	\item mathematica - notebook pro kotrolu řešení
	\item benchmark - jednoduché prográmky pro malý benchmark
	\item report - tato zpráva
\end{itemize}

\end{document}
